\chapter{DPLL}

The Davis-Putnam-Logemann-Loveland (DPLL) algorithm is a complete
backtracking based search algorithm for finding satisfiable assignments for SAT
problems. It is, dare I say, the foundational algorithm for modern SAT solvers,
as the state-of-the-art algorithm used by basically every SAT solver builds upon
the foundations that DPLL provides, as we will see later on.

\noindent It consists of three components:

\begin{itemize}
    \item Recursive backtracking
    \item Unit propagation
    \item Pure literal elimination
\end{itemize}
How does it work?

Imagine a binary tree of variables, each with two branches representing their
assignment, one for true and one for false. Suppose we don't know anything
about our CNF equation. That is, from looking at it, any variable could have
any value, and we don't know how each assignment helps us in finding out
whether or not the equation is satisfiable. How can we figure out if we can
satisfy the equation without just brute forcing it?

The answer is \sout{making an ass out of you and me} to make assumptions.

\begin{figure}[H]
    \centering
    \begin{tikzpicture}
        \tikzset{edge from parent/.style={draw,edge from parent 
            path={(\tikzparentnode.south)-- +(0,-8pt)-| (\tikzchildnode)}}}
        \Tree [.a
                    [.b
                        [.c
                            [.SAT ]
                            [.UNSAT ]
                        ]
                        [.c
                            [.SAT ]
                            [.UNSAT ]
                        ] 
                    ]
                    [.b
                        [.c
                            [.SAT ]
                            [.SAT ]
                        ]
                        [.c
                            [.UNSAT ]
                            [.UNSAT ] 
                        ] 
                    ] 
                ]
    \end{tikzpicture}
    \caption{Example DPLL decision tree (the outcomes may not correspond to an
    actual CNF instance, sue me).}
\end{figure}

Going back to our hypothetical, we don't know anything about anything. We
could have a giant decision tree to explore with an exponential number of leaves
to evaluate. How does making assumptions help us? Do we just run a random number generator 
and if it returns 1, we go ``okay, it's satisfiable,'' and false otherwise?

No. We make assumptions about our literal assignments.Unlike brute forcing where we try 
every possible combination of variable assignments at once, in DPLL we start with a literal, 
we'll call it $A$. Suppose there is nothing in the equation that indicates whether or not 
$A$ should be true or false; either could be valid.

So to tackle our problem, initially we will assume $A$ is true and see where
that leads us. We propagate that value throughout our clauses, and suppose
there is a clause that can't be satisfied. That's a contradiction, because if the
equation was satisfiable then $A = True$ couldn't have led us there.

So we've figured out that at the very least $A = True$ leads to a
contradiction. We've now eliminated half of the search space immediately and only need to
consider the cases where $A = False$.

This may not make total sense immediately, so instead of giving one explanation
of the algorithm and then building it, we will build up the algorithm and explain
how each part works.